%! Solving Radical Equations \& Extraneous Solutions — Unit 6, Lesson 6.5 — Warm-Up
\documentclass[10pt]{article}
\usepackage{algebra2-article}
\usepackage{algebra2-boxes}
\newcommand{\TallMath}[1]{$\displaystyle #1\rule[-1.4em]{0pt}{3.2em}$}

\begin{document}

\pageheader{Unit 6, Lesson 6.5}{Warm-Up}
\namedateperiod

\vspace{0.10in}

\begin{notesbox}{Getting Ready: Skills We'll Use Today}
\small Three short items. The third one is a trap, and you will meet it again in about ten minutes.

\vspace{4pt}
\textbf{1. Undoing a radical (Lessons 6.1 and 6.3).} Simplify each expression. Assume $x$ is large
enough that every radical is real.
\begin{center}
\renewcommand{\arraystretch}{1.7}
\begin{tabular}{@{}l l l@{}}
(a) \; $\left(\sqrt{x-3}\,\right)^{2}=$ \blank{2.2cm} &
(b) \; $\left(\sqrt[3]{2x}\,\right)^{3}=$ \blank{2.2cm} &
(c) \; $\left(\sqrt{x}+3\right)^{2}=$ \blank{3.4cm} \\
\end{tabular}
\end{center}
\vspace{-4pt}
Parts (a) and (c) both squared something containing $\sqrt{x}$. One of them came out clean and the
other came out \emph{worse}. What was different about the one that came out clean?
\writeline

\vspace{0.30in}

\textbf{2. A one-way street.} Answer each question \textbf{yes} or \textbf{no}.
\begin{center}
\renewcommand{\arraystretch}{1.6}
\begin{tabularx}{0.96\linewidth}{@{}X c@{}}
\hline
If you know that $x=-4$, does it follow that $x^{2}=16$? & \blank{1.6cm} \\
If you know that $x^{2}=16$, does it follow that $x=4$? & \blank{1.6cm} \\
\hline
\end{tabularx}
\end{center}
\vspace{-4pt}
Finish the sentence. \emph{Squaring both sides of a true equation always gives another true
equation, but going backwards} \blank{6.4cm}.

\vspace{4pt}
\textbf{3. Factoring, from Unit 3.} Solve by factoring.
\[
  x^{2}-11x+18=0 \qquad\Longrightarrow\qquad x=\blank{1.4cm} \ \text{ or } \ x=\blank{1.4cm}
\]
\vspace{-10pt}

\vspace{0.55in}

Hold on to those two numbers. Today an equation will hand you \emph{both} of them, and only one of
them will be telling the truth.

\end{notesbox}

\end{document}
